\chapter{Implementation}\label{ch:implementation}

\section{Introduction}\label{sec:impl-intro}

This chapter describes the implementation of the NeuroShell Intent Recognition Engine. Implementation details are derived strictly from the actual source code in \texttt{src/}, configuration in \texttt{config/}, data files in \texttt{data/}, and supporting files (\texttt{README.md}, \texttt{requirements.txt}, \texttt{pytest.ini}). The implementation follows a modular architecture with clear separation between preprocessing, inference, validation, middleware, schemas, API layer, audit, and utilities.

\section{Technology Stack}\label{sec:impl-tech}

The technology stack as implemented is summarised below (from \texttt{README.md}):

\begin{table}[htbp]
  \centering
  \caption{Technology stack of the NeuroShell IRE}
  \label{tab:tech-stack}
  \sliittablestyle
  \begin{tabular}{lp{0.32\textwidth}p{0.42\textwidth}}
    \toprule
    \textbf{Layer} & \textbf{Technology} & \textbf{Purpose} \\
    \midrule
    LLM & Gemma 4 E4B / 27B-IT, Qwen2.5-Coder 7B & Local instruction-following model for intent classification \\
    Inference Runtime & Ollama & Local model serving via REST API \\
    API Framework & FastAPI & Async HTTP API with OpenAPI/Swagger and Pydantic integration \\
    Validation & Pydantic v2 + jsonschema & Strict schema enforcement and validation \\
    Preprocessing & spaCy (\texttt{en\_core\_web\_sm}) & NLP preprocessing pipeline \\
    Logging & structlog & Structured JSON logging for auditability \\
    Rate Limiting & slowapi & Per-IP rate limiting (30 req/min on \texttt{/parse}) \\
    Testing & pytest + httpx.AsyncClient & Unit and async integration testing \\
    Containerisation & Docker + docker-compose & Reproducible deployment \\
    Fine-tuning & LoRA + PEFT + bitsandbytes & Parameter-efficient adaptation (scaffolded) \\
    \bottomrule
  \end{tabular}
\end{table}

All dependencies are specified in \texttt{requirements.txt}. The system is designed to run offline (no external API calls at inference time).

\section{Project Structure}\label{sec:impl-structure}
